\section{Asignaciones}
\subsection{Crear Asignación}

\casodeprueba{CP1.1}
{Crear asignación ok}
{Token JWT con roles suficientes (ASIGNACIONES\_CREATE), usuario existente en gestion-usuarios y vehículo existente en vehiculos-service sin asignación activa.}
{JSON con userId válido (UUID v4), vehicleId válido (UUID v4) y notas opcionales válidas (POST /asignaciones).}
{\begin{enumerate} [leftmargin=*, nosep]
    \item Código de respuesta 201
    \item Respuesta DTO con datos correctos de la asignación creada.
    \item Buscar por ID compuesto (userId y vehicleId) para confirmar su existencia.
\end{enumerate}}

\casodeprueba{CP1.2}
{Crear asignación - Ya existe una asignación activa del vehículo con el mismo usuario}
{Token JWT con roles suficientes y una asignación activa ya creada entre el usuario y el vehículo.}
{JSON con el mismo userId y vehicleId de la asignación activa existente.}
{\begin{enumerate} [leftmargin=*, nosep]
    \item Código de respuesta 409
    \item Mensaje de error indicando conflicto o que ya existe una asignación activa para el usuario y el vehículo.
\end{enumerate}}

\casodeprueba{CP1.3}
{Crear asignación - Ya existe asignación activa con otro usuario}
{Token JWT con roles suficientes y el vehículo ya tiene una asignación activa con otro usuario diferente.}
{JSON con un nuevo userId y el vehicleId que ya está asignado activamente a otro propietario.}
{\begin{enumerate} [leftmargin=*, nosep]
    \item Código de respuesta 409
    \item Mensaje de error indicando que el vehículo ya está asignado activamente a otro usuario.
\end{enumerate}}

\casodeprueba{CP1.4}
{Crear asignación - Ya existe una asignación inactiva con el mismo usuario}
{Token JWT con roles suficientes y una asignación inactiva (activa = false) previa entre el usuario y el vehículo.}
{JSON con el mismo userId y vehicleId que poseen un registro inactivo.}
{\begin{enumerate} [leftmargin=*, nosep]
    \item Código de respuesta 409
    \item Mensaje de error indicando que existe una asignación inactiva y sugiriendo usar PATCH para reactivarla.
\end{enumerate}}

\casodeprueba{CP1.5}
{Crear asignación - Usuario o vehículo no existe}
{Token JWT con roles suficientes.}
{JSON con userId o vehicleId que no existen en sus respectivos microservicios.}
{\begin{enumerate} [leftmargin=*, nosep]
    \item Código de respuesta 404
    \item Mensaje de error indicando que el usuario o el vehículo no existe.
\end{enumerate}}

\casodeprueba{CP1.6}
{Crear asignación - userId inválido (No UUID v4)}
{Token JWT con roles suficientes.}
{JSON con userId inválido (ej: '12345', 'no-uuid').}
{\begin{enumerate} [leftmargin=*, nosep]
    \item Código de respuesta 400
    \item Mensaje de error de validación para userId.
\end{enumerate}}

\casodeprueba{CP1.7}
{Crear asignación - vehicleId inválido (No UUID v4)}
{Token JWT con roles suficientes.}
{JSON con vehicleId inválido (ej: '12345', 'no-uuid').}
{\begin{enumerate} [leftmargin=*, nosep]
    \item Código de respuesta 400
    \item Mensaje de error de validación para vehicleId.
\end{enumerate}}

\casodeprueba{CP1.8}
{Crear asignación - Notas muy largas (> 500 caracteres)}
{Token JWT con roles suficientes.}
{JSON con notas de más de 500 caracteres.}
{\begin{enumerate} [leftmargin=*, nosep]
    \item Código de respuesta 400
    \item Mensaje de error de validación para notas.
\end{enumerate}}

\casodeprueba{CP2.1}
{Obtener todas las asignaciones ok}
{Token JWT con roles suficientes (ASIGNACIONES\_READ) y asignaciones creadas en el sistema.}
{Ninguna (GET /asignaciones)}
{\begin{enumerate} [leftmargin=*, nosep]
    \item Código de respuesta 200
    \item Cantidad de registros obtenidos en el array acorde al total existente.
\end{enumerate}}

\casodeprueba{CP2.2}
{Obtener flota de vehículos por propietario ok}
{Token JWT con roles suficientes (ASIGNACIONES\_READ) y vehículos asignados a un propietario específico.}
{UUID de propietario como parámetro (GET /asignaciones/propietario/:userId)}
{\begin{enumerate} [leftmargin=*, nosep]
    \item Código de respuesta 200
    \item Cantidad de registros obtenidos coincidente con la flota del propietario.
\end{enumerate}}

\casodeprueba{CP2.3}
{Obtener flota de vehículos por propietario sin resultados}
{Token JWT con roles suficientes.}
{UUID de propietario sin vehículos asignados.}
{\begin{enumerate} [leftmargin=*, nosep]
    \item Código de respuesta 200 y cantidad de registros 0 (array vacío).
\end{enumerate}}

\casodeprueba{CP2.4}
{Obtener flota de vehículos por propietario - userId inválido}
{Token JWT con roles suficientes.}
{UUID de propietario con formato inválido (GET /asignaciones/propietario/abc).}
{\begin{enumerate} [leftmargin=*, nosep]
    \item Código de respuesta 400
    \item Mensaje de error de validación (ParseUUIDPipe).
\end{enumerate}}

\casodeprueba{CP3.1}
{Buscar asignación activa por vehículo ok}
{Token JWT con roles suficientes (ASIGNACIONES\_READ) y una asignación activa asociada a un vehículo.}
{UUID del vehículo como parámetro (GET /asignaciones/vehiculo/:vehicleId/activo)}
{\begin{enumerate} [leftmargin=*, nosep]
    \item Código de respuesta 200
    \item Datos de la respuesta coinciden con la asignación activa del vehículo.
\end{enumerate}}

\casodeprueba{CP3.2}
{Buscar asignación activa por vehículo - No encontrada}
{Token JWT con roles suficientes.}
{UUID de vehículo sin asignación activa.}
{\begin{enumerate} [leftmargin=*, nosep]
    \item Código de respuesta 404
    \item Mensaje de error indicando que no existe asignación activa para el vehículo.
\end{enumerate}}

\casodeprueba{CP3.3}
{Buscar asignación activa por vehículo - vehicleId inválido}
{Token JWT con roles suficientes.}
{UUID de vehículo con formato inválido (GET /asignaciones/vehiculo/abc/activo).}
{\begin{enumerate} [leftmargin=*, nosep]
    \item Código de respuesta 400
    \item Mensaje de error de validación (ParseUUIDPipe).
\end{enumerate}}

\casodeprueba{CP3.4}
{Buscar asignación específica por clave compuesta (userId y vehicleId) ok}
{Token JWT con roles suficientes (ASIGNACIONES\_READ) y una asignación creada con un userId y vehicleId específicos.}
{UUID de usuario y UUID de vehículo como parámetros (GET /asignaciones/:userId/:vehicleId)}
{\begin{enumerate} [leftmargin=*, nosep]
    \item Código de respuesta 200
    \item Datos de la respuesta coinciden con la asignación buscada.
\end{enumerate}}

\casodeprueba{CP3.5}
{Buscar asignación específica por clave compuesta - No encontrada}
{Token JWT con roles suficientes.}
{Combinación de userId y vehicleId inexistente en el sistema.}
{\begin{enumerate} [leftmargin=*, nosep]
    \item Código de respuesta 404
    \item Mensaje de error indicando que la asignación no fue encontrada.
\end{enumerate}}

\casodeprueba{CP3.6}
{Buscar asignación específica por clave compuesta - Parámetros inválidos}
{Token JWT con roles suficientes.}
{userId o vehicleId con formato inválido (GET /asignaciones/abc/def).}
{\begin{enumerate} [leftmargin=*, nosep]
    \item Código de respuesta 400
    \item Mensaje de error de validación (ParseUUIDPipe).
\end{enumerate}}

\casodeprueba{CP4.1}
{Actualizar asignación ok}
{Token JWT con roles suficientes (ASIGNACIONES\_UPDATE) y una asignación creada.}
{DTO de actualización válido (activa: false, notas: 'Notas actualizadas') y parámetros userId y vehicleId (PATCH /asignaciones/:userId/:vehicleId).}
{\begin{enumerate} [leftmargin=*, nosep]
    \item Código de respuesta 200
    \item Datos de respuesta coinciden con los cambios aplicados.
    \item Buscar por clave compuesta para confirmar el cambio en la base de datos.
\end{enumerate}}

\casodeprueba{CP4.2}
{Actualizar asignación - Not Found}
{Token JWT con roles suficientes.}
{userId y vehicleId inexistentes y DTO de actualización válido.}
{\begin{enumerate} [leftmargin=*, nosep]
    \item Código de respuesta 404
    \item Mensaje de error de asignación no encontrada.
\end{enumerate}}

\casodeprueba{CP4.3}
{Actualizar asignación - Campo activa inválido (No booleano)}
{Token JWT con roles suficientes y una asignación creada.}
{DTO de actualización con `activa` en formato no booleano (ej: 'si' o 1).}
{\begin{enumerate} [leftmargin=*, nosep]
    \item Código de respuesta 400
    \item Mensaje de error de validación para el campo activa.
\end{enumerate}}

\casodeprueba{CP4.4}
{Actualizar asignación - Notas muy largas (> 500 caracteres)}
{Token JWT con roles suficientes y una asignación creada.}
{DTO de actualización con notas de más de 500 caracteres.}
{\begin{enumerate} [leftmargin=*, nosep]
    \item Código de respuesta 400
    \item Mensaje de error de validación para el campo notas.
\end{enumerate}}

\casodeprueba{CP4.5}
{Actualizar asignación - Parámetros inválidos (UUID incorrecto)}
{Token JWT con roles suficientes.}
{userId o vehicleId con formato inválido en la URL.}
{\begin{enumerate} [leftmargin=*, nosep]
    \item Código de respuesta 400
    \item Mensaje de error de validación (ParseUUIDPipe).
\end{enumerate}}

\casodeprueba{CP5.1}
{Eliminar asignación ok (Baja lógica)}
{Token JWT con roles suficientes (ASIGNACIONES\_DELETE) y una asignación activa creada.}
{userId y vehicleId como parámetros (DELETE /asignaciones/:userId/:vehicleId)}
{\begin{enumerate} [leftmargin=*, nosep]
    \item Código de respuesta 200
    \item Buscar por clave compuesta para confirmar que el estado activa cambió a false.
\end{enumerate}}

\casodeprueba{CP5.2}
{Eliminar asignación - Not Found}
{Token JWT con roles suficientes.}
{userId y vehicleId inexistentes.}
{\begin{enumerate} [leftmargin=*, nosep]
    \item Código de respuesta 404
    \item Mensaje de error de asignación no encontrada.
\end{enumerate}}

\casodeprueba{CP5.3}
{Eliminar asignación - Conflicto (Ya estaba inactiva)}
{Token JWT con roles suficientes y una asignación previamente dada de baja (activa = false).}
{userId y vehicleId de la asignación inactiva.}
{\begin{enumerate} [leftmargin=*, nosep]
    \item Código de respuesta 409
    \item Mensaje de error indicando que la asignación ya estaba inactiva.
\end{enumerate}}

\casodeprueba{CP5.4}
{Eliminar asignación - Parámetros inválidos (UUID incorrecto)}
{Token JWT con roles suficientes.}
{userId o vehicleId con formato inválido en la URL.}
{\begin{enumerate} [leftmargin=*, nosep]
    \item Código de respuesta 400
    \item Mensaje de error de validación (ParseUUIDPipe).
\end{enumerate}}

